\documentclass[utf8, bachelor]{gradu3}
% Jos työ on kandidaatintutkielma eikä pro gradu, käytä ylläolevan asemesta
%\documentclass[utf8,bachelor]{gradu3}
% Jos kirjoitat englanniksi, käytä ylläolevan asemesta
%\documentclass[utf8,english]{gradu3}
% tai
%\documentclass[utf8,bachelor,english]{gradu3}

\usepackage{graphicx} % kuvien mukaan ottamista varten

\usepackage{amsmath} % hyödyllinen jos tekstisi sisältää matikkaa, 

\usepackage{booktabs} % hyvä kauniiden taulukoiden tekemiseen

\usepackage{longtable}

% HUOM! Tämän tulee olla viimeinen \usepackage koko dokumentissa!
\usepackage[bookmarksopen,bookmarksnumbered,linktocpage]{hyperref}

\addbibresource{Annan-kandi.bib} % Lähdetietokannan tiedostonimi

\begin{document}

\title{Koneoppimismallien hyödyntäminen urheilijoiden rasitusvammojen ennustamisessa}
\translatedtitle{Utilizing Machine Learning Models for Predicting Overuse Injuries in Athletes}
%\studyline{Kaikki opintosuunnat}
\avainsanat{rasitusvammat, koneoppimismallit, ennustaminen}
\keywords{overuse injuries, machine learning, prediction}
\tiivistelma{%

}
\abstract{%

}

\author{Anna Eveliina Kärkkäinen}
\contactinformation{\texttt{anna.e.karkkainen@student.jyu.fi}}
% jos useita tekijöitä, anna useampi \author-komento
\supervisor{Tytti Saksa}
% jos useita ohjaajia, anna useampi \supervisor-komento

% et tarvitse tätä riviä kandidaatin- tai pro gradu -tutkielmassa!
% diplomityössä jätä komento mukaan ja muuta sulkeiden sisään "diplomityö"
\subject{Kandidaatintutkielman}
%\type{tutkimussuunnitelma} 

% vaihda tämä rivi vastaamaan tutkinto-ohjelmaasi, huom. genetiivi ja
% iso alkukirjain
\subject{Tieto- ja ohjelmistotekniikan}

\maketitle

\bigskip


\mainmatter

\chapter{Johdanto}

Koneoppimisen ala on kasvussa jatkuvasti ja sen mukana myös laskentateho paranee. Laskentatehon kasvaessa koneoppimisen 
hyödyntäminen on lisääntynyt eri tieteenaloilla, mukaan lukien urheilulääketieteessä \parencite{bhatia_exploring_2024}. Kasvava huolenaihe on 
nuorten rasitusvammojen yleistyminen, sillä nuorten urheilu on muuttunut yhä kilpailullisemmaksi \parencite{difiori_overuse_2014}. 

Tässä kandidaatintutkielmassa syvennytään aiheeseen kuinka rasitusvammat syntyvät ja mitä koneoppimisalgoritmeja käytetään niiden ennustamiseen.
Koneoppimismalleja hyödynnetään jatkuvasti erilaisissa ennustamista vaativissa tehtävissä. Tarkoituksena onkin ottaa selvää, kuinka urheilijalta 
kerättyä dataa voidaan muuntaa ja hyödyntää koneoppimismalleille sopivaksi. Tarkempana näkökulmana on juuri nuorten yleistyneet rasitusvammat ja niiden ennustaminen.

Aiheesta suoraan ei ole kovinkaan paljoa tutkimusta, joten tutkimuskysys on ajankohtainen ja mielenkiintoinen. Tutkielman tarkoituksen on siis koota erilaisia artikkeleita ja 
tutkimuksia yhteen, jotta lukija saa monipuolisen ja yksinkertaistetun kokonaisuuden rasitusvammojen syntymisestä sekä koneoppimismalleilla ennustamisesta.

Tekoälyä tässä tutkielmassa on käytetty lähteiden ymmärtämiseen ja suomentamiseen. Käytetty kielimalli on Google Gemini Pro. Muuhun, kuten omaan päättelyyn ja johtopäätöksiin, tai muuhun sisällön 
tuottamiseen ei tekoälyä ole hyödynnetty.

\chapter{Käsitteet}

\section{Rasitusvammat}

Tässä tutkielmassa puhutaan sekä urheiluvammoista että rasitusvammoista. Urheiluvammoilla tarkoitetaan sekä akuutteja että rasitusvammoja.

Rasitusvammat syntyvät ajan myötä, toisin kuin akuutit vammat, jotka syntyvät suoraan tapahtumasta. Rasitusvammoja voi syntyä, kun yksilö altistuu 
toistuvalle rasitukselle. Kun kehon kudokset kokee liiallista rasitusta, liian vähäistä palautumista tai molempia, kehon kyky palautua ylittyy \parencite{johnson_overuse_2008}.
Tarkemmin rasitusvammat ovat seurausta luuhun, lihakseen ja/tai jänteeseen kohdistuvasta kumulatiivisesta 
mikrotraumasta \parencite{brenner_overuse_2024}.

\textcite{johnson_overuse_2008} toteaa rasitusvammojen syyt jakautuvan yleensä kahteen osaan. Sisäisiin tekijöihin eli yksilöstä johtuviin ja ulkoisiin 
tekijöihin eli ympäristöstä johtuviin. Hän mainitsee sisäisiä tekijöitä olevan muun muassa anatomiset virheasennot, heikko fyysinen kunto 
sekä aikaisemmat vammat. Ulkoisia tekijöitä puolestaan ovat esimerkiksi väärät harjoittelumetodit, huonot tekniikat sekä 
valmentajien tai harjoittelu kavereiden aiheuttama paine. Ulkoisiin tekijöihin yksilön on mahdollista vaikuttaa, toisin kuin sisäisiin tekijöihin.

Nuorten kilpailullinen urheilu on muuttunut yhä vaativammaksi, mikä on johtanut rasitusvammojen yleistymiseen \parencite{weerasinghe_computer_2016}. Treenit ovat usein intensiivisiä, 
ja nuoret urheilijat saattavat altistua liialliselle rasitukselle.

\section{Koneoppimismallit}

Koneoppiminen (engl. Machine Learning) tarkoittaa automatisoitua prosessia, jossa datasta tunnistetaan säännönmukaisuuksia 
tehtävien, kuten luokittelun ja ennustamisen, suorittamiseksi \parencite{black_introduction_2023}. Koneoppimismallit jaetaan 
yleisesti kolmeen eri kategoriaan: ohjattu oppiminen, ohjaamaton oppiminen ja vahvistettu oppiminen. 

Tässä työssä käsiteltävät menetelmät rajautuvat ohjattuun oppimiseen, sillä ohjattun oppimisen menetelmät ovat käytetyimpiä ennsutamiseen \parencite{singh_machine_2021}.

\subsection{Ohjattu oppiminen}

Ohjatussa oppimisessa mallille syötetään dataa, jonka lopputulokset tunnetaan. Malli oppii tunnistamaan datasta säännönmukaisuuksia ja oppii näin luokittelemaan tai ennustamaan 
havaitsemattomia lopputuloksia uudesta datasta \parencite{black_introduction_2023}. Ohjatun oppimisen menetelmiä on useita, joista tässä työssä merkittävimpiä ovat seuraavat:

\begin{longtable}{| p{.35\textwidth} | p{.65\textwidth} |}
\hline
\textbf{Menetelmä} & \textbf{Selitys} \\ \hline
Extreme Gradient Boosting (XGBoost) & Koostuu gradienttitehostetuista päätöspuista. Tehostamisen perusideana on rakentaa peräkkäin alipuita alkuperäisestä puusta siten, että jokainen seuraava puu vähentää edellisen puun virheitä \parencite{8906396}. \\ \hline
Lineaarinen ja logistinen regressio (engl. Linear and logistic regression) & Lineaarinen regressio etsii suoran, joka sopii dataan parhaiten tietyn matemaattisen kriteerin mukaisesti. Logistinen regressio etsii puolestaan dataan parhaiten sopivan viivan sen jälkeen, kun siihen on sovellettu logistista funktiota \parencite{black_introduction_2023}. \\ \hline
Lähimmän naapurin menetelmä (engl. K-Nearest Neighbors, K-NN) & Algoritmia käytetään sekä regressio- että luokitteluongelmien ratkaisemiseen. Käytetään ennustamiseen samankaltaisten tapausten keskiarvon ja mediaanin perusteella sekä laskemaan tulosdata perustuen K-määrän samankaltaisimpien tapausten esiintymistiheyteen \parencite{singh_machine_2021}. \\ \hline
Malliyhdistelmät (engl. Ensembles) & Yhdistää useita koneoppimismalleja tehdäkseen luokittelua tai ennustamista. Uudet luokittelut ja ennusteet perustuvat kunkin osamallin antamien ennusteiden yhteisymmärrykseen \parencite{black_introduction_2023}. \\ \hline
Naiivi Bayes (engl. Naive Bayes) & Perustuu Bayesin teoreemaan ja oletukseen siitä, että mallin kunkin piirteen esiintyminen on riippumatonta muista \parencite{singh_machine_2021}. Se on luokittelija, joka muokkaa alkuperäistä ennakkotodennäköisyyttä sen perusteella, kuinka usein tiettyjä piirteitä havaitaan aineiston tapauksissa \parencite{black_introduction_2023}.\\ \hline
Neuroverkot (engl. Neural networks) & Järjestetty, toisiinsa kytketty, solmupisteiden kokoelma, joka on jaettuna useisiin kerroksiin. Jokainen solmu vastaanottaa signaalin ja muuntaa sen tulosteeksi. Solmuille määritetään painokertoimet, jotka määräävät, miten ne vaikuttavat lopulliseen tulokseen \parencite{black_introduction_2023}. \\ \hline
Päätöspuu (engl. Decision trees) & Data jaetaan yhä samankaltaisempiin ryhmiin perustuen muuttujaan, joka maksimoi tuloksena olevien ryhmien samankaltaisuuden. Uudet luokittelut ja ennusteet tehdään arvioimalla uutta yksilöä näiden jakokriteerien mukaisesti \parencite{black_introduction_2023}.\\ \hline
Satunnaismetsä (engl. Random forest) & Päätöspuihin perustuva ryhmämenetelmä, joka yhdistää useita itsenäisesti satunnaistettuja päätöspuita. Jokainen puu rakennetaan hyödyntämällä satunnaista otantaa ja satunnaisia muuttujajoukoja, mikä varmistaa mallin vakauden ja vähentää virhealttiutta \parencite{breiman_random_2001}. \\ \hline
Tukivektorikoneet (engl. Support vector machines, SVM) & Malli jakaa havaintoaineiston osiin n-ulotteisen hypertason perusteella. Uudet yksilöt luokitellaan sen mukaan, minkä rajatun alueen sisälle heidän ominaisuutensa osuvat \parencite{black_introduction_2023}.\\ \hline
    \caption{Ohjatun oppimisen menetelmiä} \label{long}\\
\end{longtable}


\chapter{Rasitusvammojen ennustaminen koneoppimismalleilla}

Edellisessä luvussa esiteltiin erilaisia ohjatun oppimisen menetelmiä, joita voidaan hyödyntää ennustamiseen. Menetelmän valitseminen riippuu ennustettavasta ongelmasta, 
datan luonteesta ja ennustamisen vaatimuksista. 

Useat tutkimukset käsittelevät urheiluvammoja yleisellä tasolla, eikä niissä eritellä tarkemmin, onko kyseessä akuutti vai rasitusvamma.
\textcite{murugan_comparison_2025} tutkimuksessa parhaat tulokset urheiluvammojen ennustamiseen saavutettiin käyttämällä satunnaismetsää ja tukivektorikoneita. 
\textcite{leckey_machine_2025} systemaattisessa kirjallisuuskatsauksessa parhaiksi algoritmeiksi puolestaan osoittautuivat satunnaismetsät sekä Extreme Gradient Boosting.

\section{Datan keruu ja sen esikäsittely}

Datan esikäsittely on tärkeää, sillä koneoppimismallit vaativat datan olevan tietyssä muodossa. Monesti urheilijalta kerätty data on esimerkiksi saatua 
dataa urheilukellosta tai kyselyiden pohjalta saatua dataa. Tästä johtuen dataa on esikäsiteltävä, jotta malli osaa sitä hyödyntää. Saatua dataa muun muassa puhdistetaan, 
normalisoidaan ja muunnetaan eri muotoon \parencite{murugan_comparison_2025}. Esikäsittelyyn kuuluu myös esimerkiksi datan imputointi eli paikkaus puuttuvien arvojen 
hallitsemiseksi \parencite{yuan_machine_2026}. 

Rasitusvammoja tutkiessa yleisin datan keruu muoto on erilaiset kyselyt, joihin urheilijat vastaavat. Tutkimuksessa \parencite{murugan_comparison_2025} 
vertailtiin koneoppimismalleja urheilijoiden loukkaantumisten ennustamiseen. Tutkimuksessa dataa kerättiin useilta urheilijoilta määritellyn ajanjakson aikana. 
Jokainen aineiston merkintä vastasi yhtä harjoitus päivää ja sisälsi muun muassa seuraavat ominaisuudet: nr.sessions (urheilijan yhden päivän aikana suorittamien 
harjoitusten kokonaismäärä), total km (kaikkien harjoitusten aikana kertynyt kokonaiskilometrimäärä), km z3-4 (sykealueilla 3 ja 4 taitettu matka) sekä strength 
training (päivän aikana tehty voimaharjoittelun kokonaismäärä). 

\textcite{de_leeuw_personalized_2022} tutkimuksessa puolestaan urheilijat vastasivat päivittäin OSTRC-kyselyyn (engl. Oslo Sports Trauma 
Research Center questionnaire). Tutkimuksen kysymykset jakautuivat neljään kategoriaan: 1) Vaikeus osallistua normaaliin harjoitteluun ja kilpailuun, 2) Vähentynyt harjoittelumäärä,
3) Vaikutus suorituskykyyn ja 4) Koetut oireet/vaivat. Datan analysoimisessa nämä neljä kysymyskategoriaa toimivat vastemuuttujina, joita analysoitiin yksi kerrallaan.
Ennustemuuttujina käytettiin sisäistä ja ulkoista harjoituskuormitusta sekä koettua hyvinvointia.


\section{Rasitusvammojen ennustaminen}

Koneoppimismallien viimeaikainen kehittyminen on avannut uusia ovia monimutkaisten tietoaineistojen analysointiin, mikä mahdollistaa kaavojen löytämisen, 
joita perinteiset menetelmät saattavat ohittaa \parencite{murugan_comparison_2025}. Nykyisiä koneoppimismenetelmiä voidaan käyttää urheiluvammojen riskitekijöiden 
havaitsemiseen sekä tunnistamaan urheilijoita, joilla on korkea vammautumisriski \parencite{van_eetvelde_machine_2021}.

\textcite{de_leeuw_personalized_2022} tutkimus esitteli Subgroup discovery -algoritmin rasitusvammojen ennustamiseen. Tutkimuksessa esitelty koneoppimisalgoritmi on ohjatun oppimisen 
menetelmä, mikä analysoi yksittäisten tai vain muutamien selittävien muuttujien vaikutusta valittuun vastemuuttujaan ja nostaa esiin ainoastaan merkittävimmät selittävät muuttujat. 
Tämä algoritmi mahdollistaa yksilöllisten yhteyksien havainnoinnin sekä saadut tulokset ovat suoraviivaisia ja helposti yhdistettävissä käytännön treeneihin.

\textcite{InjuryPredictioninCompetitiveRunnersWithMachineLearning} tutkimuksessaan puolestaan esittelivät lähestymistapana Extreme Gradient Boosting -algoritmin ennustamaan urheilijoiden 
loukkaantumisia. Tutkimuksessa kerrottiin, kuinka akuutit vammat ovat hankalasti ennustettavissa, toisin kuin rasitusvammat. 

Muissa tutkimuksissa ja artikkeleissa, kuten \textcite{murugan_comparison_2025}, \textcite{bhatia_exploring_2024} sekä \textcite{yuan_machine_2026} urheiluvammojen ennustamisesta 
puhuttiin yleisellä tasolla, mutta ei spesifioitu tarkemmin oliko kyseessä akuutit- vai rasitusvammat. Muun muassa satunnaismetsä, logistinen regressio, Extreme Gradient Boosting 
ja lähimmän naapurin menetelmä mainittiin toimivina algoritmeina.

\section{Lopputulokset}

Urheiluvammojen ennustamista koneoppimismalleilla on tutkittu lähivuosina enemmän ja enemmän. Kuitenkin rasitusvammojen tutkiminen on jäänyt vähäiseksi. Rsitusvammat 
syntyvät ajan myötä, tietynlaisesta toistuvasta rasituksesta. Koneoppimismallit puolestaan tunnistavat datasta toistuvia kuvioita ja oppivat näin ennustamaan uudesta datasta 
lopputuloksia. Näin ollen koneoppimismallien hyödyntämistä rasitusvammojen ennustamiseen tulisi tutkia enemmän tulevaisuudessa, sillä kuten alussa mainittiin nuorten rasitusvammojen 
yleistyminen on kasvava huolenaihe.

Lopputuloksena voidaan todeta, että koneoppimismallit auttavat tunnistamaan yksilöllisiä riskitekijöitä, jotka johtavat rasitusvammoihin. Näiden riskitekijöiden tunnistaminen 
mahdollistaa ennaltaehkäisemään jo alkuvaiheessa olevia rasitusvammoja.

\chapter{Yhteenveto}

Koneoppimismalleja käytetään jatkuvasti erilaisissa ennustamisissa. Tässä tutkielmassa paneuduttiin siihen, kuinka urheilijalta kerättyä 
dataa voidaan muuntaa ja hyödyntää koneoppimismalleille sopivaksi.

Lopussa yhteenvedossa tulee ilmi tiivis kokonaisuus juuri kertomistani asioista. Yhteenvedossa tulee näkymään selkeä ymmärrykseni aiheesta.

\chapter{Johtopäätökset}

Tutkielaman aiheesta ei kovin paljoa tutkimusta löydy. Aihe on kuitenkin hyvin tuore ja varmasti lisää tutkimusta tulevaisuudessa tehdään. 
Koneopimismallit kehittyvät jatkuvasti ja sitä myöten tutkimukset siitä kuinka niitä voidaan ennustamiseen hyödyntää, lisääntyvät.

KYSYMYS: Tuntuu, että olisi fikusa olla yhteenveto ja johtopäätökset, mutta jotenkin jumitun siihen, että miten ne tulisi koota. Onko tähän jotain hyviä konkreettisia vinkkejä tai tulisiko ne vaan yhdistää yhdeksi?

\printbibliography


\end{document}
